\documentclass[border=10pt]{standalone}
\usepackage{fontspec}
\setmainfont{Noto Serif}
\newfontfamily\bengalifont{Noto Serif Bengali}[Script=Bengali]
\usepackage[Latin,Bengali]{ucharclasses}
\setTransitionsFor{Bengali}{\begingroup\bengalifont}{\endgroup}
\usepackage{newunicodechar}
\newunicodechar{।}{{\bengalifont\char"0964}}
\newunicodechar{॥}{{\bengalifont\char"0965}}

\usepackage[table]{xcolor}
\usepackage{booktabs,tabularx,array}
\definecolor{bapilrule}{HTML}{1F2937}
\definecolor{bapilstripe}{HTML}{F3F4F6}
\renewcommand{\arraystretch}{1.28}
\setlength{\tabcolsep}{5pt}

\begin{document}
\begin{minipage}{17cm}
{\large\bfseries সারণি ৫৬।}\ {\large গবেষণা-কর্মসূচি ও পরিকল্পিত জ্ঞান-উৎপাদ}\par\vspace{7pt}
\small
\rowcolors{2}{bapilstripe}{white}
\begin{tabularx}{\linewidth}{@{}>{\hsize=0.372\hsize\raggedright\arraybackslash}X>{\hsize=1.596\hsize\raggedright\arraybackslash}X>{\hsize=1.170\hsize\raggedright\arraybackslash}X>{\hsize=1.277\hsize\raggedright\arraybackslash}X>{\hsize=0.585\hsize\raggedright\arraybackslash}X@{}}
\arrayrulecolor{bapilrule}\toprule
\rowcolor{white}\textbf{আইডি} & \textbf{গবেষণা প্রশ্ন} & \textbf{পদ্ধতি} & \textbf{পরিকল্পিত উৎপাদ} & \textbf{সময়} \\
\midrule
RQ1 & বাংলা ও ইংরেজি নীতি-টেক্সটে শব্দভিত্তিক, ঘন ও মিশ্র উদ্ধার কতটা কার্যকর, এবং দুই ভাষার ব্যবধান কীসে কমে? & BanglaPolicyQA-তে বেঞ্চমার্ক পরীক্ষা; প্রমিতকরণ ও শব্দকোষ সম্প্রসারণের প্রভাব বিশ্লেষণ & বেঞ্চমার্ক ও পদ্ধতি-প্রবন্ধ; উন্মুক্ত কোড ও ডেটা & ২০২৭ থেকে ২০২৮ \\
RQ2 & বিরতি ও উদ্ধৃতি-যাচাই কি এমনভাবে নির্ধারণ করা যায় যাতে কর্মকর্তারা কম অসমর্থিত দাবি দেখেন অথচ উপকারী উত্তর না হারান? & সীমা-নির্ধারণ সমীক্ষা; বিশ্বস্ততার মানবিক বিচার & মূল্যায়ন-প্রবন্ধ; নির্ধারণ-নির্দেশিকা & ২০২৭ থেকে ২০২৯ \\
RQ3 & উদ্ধৃতিযুক্ত অনুসন্ধান-সরঞ্জাম কি নীতি-বিশ্লেষণের গতি ও মান বদলায়? & পাইলট ব্যবহারকারীদের সঙ্গে মিলিত-কাজ সমীক্ষা; সাক্ষাৎকার & প্রমাণ-প্রাপ্তির সময় ও মানের অভিজ্ঞতাভিত্তিক সমীক্ষা & ২০২৮ থেকে ২০২৯ \\
RQ4 & নতুন ডেটা আইনের অধীনে একটি জনপ্রশাসনে ঝুঁকি-স্তর, গেট ও স্বচ্ছতা নিবন্ধন কীভাবে কার্যকর করা যায়? & ডিজাইন সায়েন্স ও BAPIL শাসনের কেস স্টাডি & ব্লুপ্রিন্টের কাঠামো বিষয়ে পিয়ার-রিভিউড জার্নাল প্রবন্ধ; নীতি-সংক্ষেপ & পাণ্ডুলিপি ২০২৭ \\
RQ5 & OCR-এর মান ও ইউনিকোড প্রমিতকরণ বাংলায় পরবর্তী উদ্ধারকে কীভাবে প্রভাবিত করে? & স্ক্যানকৃত ও জন্মগতভাবে ডিজিটাল নথিতে নিয়ন্ত্রিত পরীক্ষা & ডেটাশিটসহ সংশোধিত টেক্সটের ডেটাসেট; সংক্ষিপ্ত প্রবন্ধ & ২০২৮ \\
RQ6 & নিম্ন-মধ্যম আয়ের দেশে কোন প্রাতিষ্ঠানিক রূপ জনস্বার্থমূলক এআই নীতি ল্যাবকে টেকসই রাখে? & তুলনামূলক প্রাতিষ্ঠানিক বিশ্লেষণ; মূল্যায়নের প্রমাণ & ব্লুপ্রিন্ট ২.০-এর জন্য প্রাতিষ্ঠানিক পর্যালোচনা & ২০২৯ \\
\bottomrule
\end{tabularx}
\par\vspace{5pt}{\footnotesize\itshape উৎপাদগুলো DOI-সহ উন্মুক্ত প্রবেশাধিকারে প্রকাশিত হবে। জার্নাল নির্বাচন পরিধি, উন্মুক্ত প্রবেশাধিকারের সুযোগ ও পিয়ার রিভিউয়ের মান অনুযায়ী; পিয়ার রিভিউয়ের কোনো ফলাফল পূর্বানুমান করা হয় না।}
\par\vspace{3pt}{\scriptsize\color{gray} উৎস: BAPIL কারিগরি ব্লুপ্রিন্ট, সংস্করণ ১.০ (২০২৬)। CC BY 4.0।}
\end{minipage}
\end{document}
